\chapter[Regularność Sobolewa rozwiązań eliptycznych]{Regularność Sobolewa rozwiązań eliptycznych: ilorazy różnicowe, oszacowania wewnętrzne i brzegowe}
\label{chap:pde15}

W rozdziale~\ref{chap:pde13} skonstruowaliśmy rozwiązania słabe
$u\in H^1$. Teraz zbadamy, kiedy równanie pozwala odzyskać drugie
pochodne $u$ w $L^2$. Zasadnicza trudność jest logiczna: nie wolno
różniczkować równania, zanim nie udowodnimy, że odpowiednie pochodne
istnieją. Ominiemy ją dzięki ilorazom różnicowym, których zastosowanie
wymaga jedynie $H^1$. W końcowej części zobaczymy, dlaczego gładkość
współczynników i brzegu naprawdę jest potrzebna.

W całym rozdziale $d\ge1$, a obszarem w $\R^d$ nazywamy niepusty,
otwarty i spójny podzbiór $\R^d$. Dla otwartych zbiorów
$U,V$ zapis $U\Subset V$ oznacza, że $\overline U$ jest zwartym
podzbiorem $V$. Symbole $H^k(U)=W^{k,2}(U)$ oraz
$\|\cdot\|_{2,U}=\|\cdot\|_{L^2(U)}$ mają znaczenie ustalone
w rozdziale~\ref{chap:pde-sobolev}. Całki po $U$ zapisujemy względem
miary Lebesgue'a $\mathrm d\lambda_d(x)$.

\section{Eliptyczność a możliwość różniczkowania rozwiązań}
\label{sec:pde15-model}
Rozważamy równanie w postaci dywergencyjnej
\begin{equation}\label{eq:pde15-basic}
 -\diver(A\nabla u)=f\qquad\text{w }\Omega,
\end{equation}
gdzie $\Omega\subset\R^d$ jest obszarem, $A=(a_{ij})$ ma współczynniki
rzeczywiste, a $f\in L^2_{\mathrm{loc}}(\Omega)$.
W całym rozdziale funkcje $u,f$ i funkcje testowe są rzeczywiste.

\begin{definicja}[Jednostajna eliptyczność i rozwiązanie słabe]
\label{def:pde15-elliptic}
Mówimy, że $A\in L^\infty_{\mathrm{loc}}(\Omega;\R^{d\times d})$
jest jednostajnie eliptyczna i ograniczona na mierzalnym
$U\subset\Omega$, jeżeli istnieją $\lambda,\Lambda>0$
oraz zbiór $E\subset U$ pełnej miary w $U$ takie, że
dla każdego $x\in E$ i każdego $\xi\in\R^d$ zachodzi
$\xi^{\mathsf T}A(x)\xi\ge\lambda|\xi|^2$ oraz
$|A(x)\xi|\le\Lambda|\xi|$. Nie wymagamy symetrii $A$.
Funkcja $u\in H^1_{\mathrm{loc}}(\Omega)$ jest rozwiązaniem słabym
równania~\eqref{eq:pde15-basic}, jeżeli dla każdego
$\varphi\in C_c^\infty(\Omega)$ zachodzi
\[
 \int_\Omega A\nabla u\cdot\nabla\varphi\,\mathrm d\lambda_d(x)
 =\int_\Omega f\varphi\,\mathrm d\lambda_d(x).
\]
Równość rozciąga się na funkcje $\varphi\in H^1(\Omega)$
o zwartym nośniku $K\subset\Omega$. Istotnie, wybieramy
otwarty $W\Subset\Omega$ zawierający $K$ i funkcję
$\psi\in C_c^\infty(W)$ równą $1$ w otoczeniu $K$.
Z lematu~\ref{lem:pde9-compact-product} dla $\psi\varphi=\varphi$
otrzymujemy ciąg $\varphi_n\in C_c^\infty(W)$ zbieżny do
$\varphi$ w $H^1(W)$. Ponieważ $A$ jest ograniczona na $W$,
a $\nabla u,f\in L^2(W)$, mamy
\[
 \left|\int_W A\nabla u\cdot\nabla(\varphi_n-\varphi)\,
 \mathrm d\lambda_d(x)\right|
 \le \|A\|_{\infty,W}\|\nabla u\|_{2,W}
       \|\nabla(\varphi_n-\varphi)\|_{2,W},
\]
\[
 \left|\int_W f(\varphi_n-\varphi)\,\mathrm d\lambda_d(x)\right|
 \le \|f\|_{2,W}\|\varphi_n-\varphi\|_{2,W}.
\]
Obie prawe strony dążą do zera, co uzasadnia przejście
do granicy w równaniu słabym.
\end{definicja}

\begin{uwaga}[Dlaczego samo $A\in L^\infty$ nie wystarcza?]
Jeśli $u$ ma tylko jedną słabą pochodną, to iloczyn $A\nabla u$
należy jedynie do $L^2_{\mathrm{loc}}$. Wyrażenie
$\diver(A\nabla u)$ jest jednak poprawnie określoną dystrybucją.
Przejście do $\nabla^2u\in L^2$ wymaga dodatkowej kontroli zmian
współczynników $A$. W poniższym twierdzeniu wystarczy warunek
$A\in W^{1,\infty}_{\mathrm{loc}}$.
\end{uwaga}

\section{Ilorazy różnicowe i ich podstawowe własności}
\label{sec:pde15-differences}
Dla $1\le j\le d$ i $h\ne0$ wprowadzamy
$\delta_h^jv(x)=\bigl(v(x+he_j)-v(x)\bigr)/h$, o ile oba punkty
należą do zbioru, na którym określona jest $v$.

\begin{lemat}[Oszacowania i całkowanie przez części]
\label{lem:pde15-difference}
Niech $U\subset\R^d$ będzie otwarty. Dla $h\ne0$ i ustalonego
$j$ oznaczamy $U_h=U\cap(U-he_j)$.
Jeżeli $v\in H^1(U)$, $K\subset U$ jest niepusty i zwarty oraz
$|h|<\operatorname{dist}(K,\R^d\setminus U)$
(przy $U=\R^d$ przyjmujemy odległość równą $+\infty$),
to $\|\delta_h^jv\|_{L^2(K)}\le\|\partial_jv\|_{L^2(U)}$.
Jeżeli $v\in L^2_{\mathrm{loc}}(U)$, $w\in L^2(U)$ ma
zwarty nośnik $K\subset U$ i $K+he_j\subset U$, to
\[
 \int_{U_h}(\delta_h^jv)w\,\mathrm d\lambda_d(x)
 =-\int_U v(\delta_{-h}^jw)\,\mathrm d\lambda_d(x),
\]
gdzie po prawej $w$ przedłużamy przez zero poza $U$.
Nie wymagamy zwartego nośnika $v$. Dla
$v\in H^1_{\mathrm{loc}}(U)$ zachodzi ponadto
$\delta_h^jv\in H^1_{\mathrm{loc}}(U_h)$ i
$\partial_k(\delta_h^jv)=\delta_h^j(\partial_kv)$
w $U_h$ dla $1\le k\le d$.
\end{lemat}
\begin{proof}
Dla funkcji gładkiej podstawowe twierdzenie rachunku całkowego
wzdłuż odcinka (również dla $h<0$) daje
$\delta_h^jv(x)=\int_0^1\partial_jv(x+the_j)\,
\mathrm d\lambda_1(t)$. Każdy przesunięty zbiór
$K+the_j$ jest zawarty w $U$. Z nierówności Cauchy'ego--Schwarza,
twierdzenia Tonellego i niezmienniczości miary Lebesgue'a
względem translacji otrzymujemy
\[
 \int_K|\delta_h^jv(x)|^2\,\mathrm d\lambda_d(x)
 \le\int_0^1\int_{K+the_j}|\partial_jv(y)|^2\,
 \mathrm d\lambda_d(y)\,\mathrm d\lambda_1(t)
 \le\int_U|\partial_jv(y)|^2\,\mathrm d\lambda_d(y).
\]
Dla ogólnej funkcji Sobolewa zbiór
$S=\{x+the_j:x\in K,\ 0\le t\le1\}$ jest zwarty jako obraz
$K\times[0,1]$ przez odwzorowanie ciągłe i zawarty w $U$.
Wybieramy otwarty $W$ taki, że $S\subset W\Subset U$.
Z twierdzenia~\ref{thm:pde9-mollify} otrzymujemy funkcje
gładkie $v_n$ na $W$, zbieżne do $v$ w $H^1(W)$.
Dla ustalonego $h$ mamy
\[
 \|\delta_h^j(v_n-v)\|_{2,K}
 \le |h|^{-1}(\|v_n-v\|_{2,K+he_j}
                  +\|v_n-v\|_{2,K})
 \le 2|h|^{-1}\|v_n-v\|_{2,W}\longrightarrow0.
\]
Stosując poprzedni rachunek do $v_n$ na $W$, a następnie
przechodząc do granicy, dostajemy
$\|\delta_h^jv\|_{2,K}\le\|\partial_jv\|_{2,W}
\le\|\partial_jv\|_{2,U}$.

W drugiej tożsamości wszystkie całki istnieją z nierówności
Cauchy'ego--Schwarza, ponieważ $v$ jest lokalnie kwadratowo
całkowalna, a $K$ oraz $K+he_j$ są zwarte w $U$.
Mamy $K\subset U_h$. Po przedłużeniu $w$ przez zero
podstawienie $y=x+he_j$ daje
\[
 \int_{U_h}(\delta_h^jv)w\,\mathrm d\lambda_d(x)
 =\frac1h\int_U v(y)\bigl(w(y-he_j)-w(y)\bigr)
 \,\mathrm d\lambda_d(y)
 =-\int_U v\,\delta_{-h}^jw\,\mathrm d\lambda_d.
\]
Wreszcie dla $\zeta\in C_c^\infty(U_h)$, po przedłużeniu
jej przez zero, z tej samej tożsamości otrzymujemy
\begin{align*}
 \int_{U_h}\delta_h^jv\,\partial_k\zeta\,\mathrm d\lambda_d
 &=-\int_U v\,\partial_k(\delta_{-h}^j\zeta)\,\mathrm d\lambda_d\\
 &=\int_U \partial_kv\,\delta_{-h}^j\zeta\,\mathrm d\lambda_d\\
 &=-\int_{U_h}\delta_h^j(\partial_kv)\zeta\,\mathrm d\lambda_d.
\end{align*}
Na każdym zwartym podzbiorze $U_h$ obie translacje
$v$ i $\partial_kv$ są kwadratowo całkowalne.
Ostatnia tożsamość daje więc postulowaną słabą pochodną.
\end{proof}

\begin{lemat}[Kryterium istnienia słabej pochodnej]
\label{lem:pde15-criterion}
Niech $U\subset\R^d$ będzie otwarty i $v\in L^2_{\mathrm{loc}}(U)$.
Jeśli dla każdego niepustego zwartego $K\subset U$ istnieją
$C_K,h_K>0$ takie, że $K+he_j\subset U$ oraz
$\|\delta_h^jv\|_{L^2(K)}\le C_K$ dla $0<|h|<h_K$, to
$\partial_jv\in L^2_{\mathrm{loc}}(U)$ i
$\|\partial_jv\|_{L^2(K)}\le C_{K'}$ dla zwartych
$K,K'\subset U$ spełniających $K\subset\operatorname{int}K'$.
\end{lemat}
\begin{proof}
Ustalmy takie $K,K'$. Wybieramy otwarty $V$ z
$K\subset V\subset\overline V\subset\operatorname{int}K'$;
jest to możliwe dzięki zwartości $K$ i otwartości
$\operatorname{int}K'$. Dla dowolnego ciągu niezerowych
$h_n\to0$ z $|h_n|<h_{K'}$ mamy
$\|\delta_{h_n}^jv\|_{2,V}\le C_{K'}$.
Z twierdzenia o słabej ciągowej zwartości ograniczonych
ciągów w przestrzeni Hilberta $L^2(V)$ wybieramy podciąg
(oznaczany tak samo), dla którego
$\delta_{h_n}^jv\rightharpoonup w$ w $L^2(V)$.
Dla $\varphi\in C_c^\infty(V)$ stosujemy poprzedni lemat
dla dostatecznie małych $|h_n|$:
\[
 \int_V\delta_{h_n}^jv\,\varphi\,\mathrm d\lambda_d
 =-\int_U v\,\delta_{-h_n}^j\varphi\,\mathrm d\lambda_d .
\]
Nośniki $\delta_{-h_n}^j\varphi$ dla dużych $n$
leżą w jednym zwartym podzbiorze $V$.
Z podstawowego twierdzenia rachunku całkowego i jednostajnej
ciągłości $\partial_j\varphi$ na $\R^d$ wynika jednostajna
zbieżność $\delta_{-h_n}^j\varphi\to\partial_j\varphi$,
a z ograniczenia nośników także zbieżność w $L^2(V)$.
Z nierówności Cauchy'ego--Schwarza prawa strona zbiega do
$-\int_V v\partial_j\varphi\,\mathrm d\lambda_d$;
lewa zbiega do $\int_V w\varphi\,\mathrm d\lambda_d$
na mocy słabej zbieżności. Stąd $w$ reprezentuje
$\partial_jv$ na $V$. Półciągłość normy przy słabej zbieżności
daje
\[
 \|w\|_{2,K}\le\|w\|_{2,V}
 \le\liminf_n\|\delta_{h_n}^jv\|_{2,V}\le C_{K'}.
\]
Każdy punkt $U$ ma takie otoczenie $V$. Na przecięciach
reprezentanty są równe prawie wszędzie z jednoznaczności
dystrybucji regularnych (stwierdzenie~\ref{prop:pde7-injective}).
Wybierając przeliczalne podpokrycie tych otoczeń (przestrzeń
euklidesowa ma przeliczalną bazę), otrzymujemy jeden
mierzalny reprezentant $\partial_jv$, lokalnie kwadratowo
całkowalny. Istotnie, każdy zwarty podzbiór $U$ pokrywają
skończenie liczne z tych otoczeń i na każdym z nich
odpowiednia całka kwadratu jest skończona.
\end{proof}

\section[Wewnętrzna regularność Sobolewa drugiego rzędu]{Wewnętrzna regularność $H^2$}
\label{sec:pde15-h2}
\begin{twierdzenie}[Wewnętrzna regularność drugiego rzędu]
\label{thm:pde15-h2}
Niech $\Omega\subset\R^d$ będzie obszarem,
$A\in W^{1,\infty}_{\mathrm{loc}}(\Omega;\R^{d\times d})$ będzie
jednostajnie eliptyczna i ograniczona na zwartych podzbiorach,
$f\in L^2_{\mathrm{loc}}(\Omega)$, a $u\in H^1_{\mathrm{loc}}(\Omega)$
rozwiązaniem słabym równania~\eqref{eq:pde15-basic}.
Wówczas $u\in H^2_{\mathrm{loc}}(\Omega)$. Dla każdych
niepustych otwartych zbiorów $U,V$ takich, że
$U\Subset V\Subset\Omega$, zachodzi
\begin{equation}\label{eq:pde15-h2-est}
 \|u\|_{H^2(U)}\le C\bigl(\|f\|_{L^2(V)}
 +\|u\|_{H^1(V)}\bigr),
\end{equation}
gdzie $C$ zależy od $U,V,d$, stałych eliptyczności na $V$
oraz $\|A\|_{W^{1,\infty}(V)}$.
\end{twierdzenie}
\begin{proof}
\emph{Krok 1: dziedziny przesunięć i równanie różnicowe.}
Wybieramy $\eta\in C_c^\infty(V)$ taką, że
$0\le\eta\le1$ i $\eta=1$ w otoczeniu $\overline U$.
Niech $K=\operatorname{supp}\eta$ i wybierzmy $\rho>0$ tak,
aby $K+\overline{B_{2\rho}(0)}\subset V$.
Ustalamy $j\in\{1,\ldots,d\}$ i $0<|h|<\rho$.
Oznaczamy $V_h=V\cap(V-he_j)$ i $w_h=\delta_h^ju$ na $V_h$.
Z lematu~\ref{lem:pde15-difference} mamy
$w_h\in H^1_{\mathrm{loc}}(V_h)$ oraz
$\nabla w_h=\delta_h^j\nabla u$.
W szczególności $K\Subset V_h$.

Dla $\varphi\in C_c^\infty(V_h)$ funkcja
$-\delta_{-h}^j\varphi$ (po przedłużeniu $\varphi$
przez zero) jest gładkim testem o zwartym nośniku w $V$.
Podstawiamy ją do równania słabego
i stosujemy dyskretne całkowanie przez części
z lematu~\ref{lem:pde15-difference}, osobno do każdej
składowej pola $A\nabla u$ oraz do $f$.
Otrzymujemy
\[
 \int_{V_h}\delta_h^j(A\nabla u)\cdot\nabla\varphi\,
 \mathrm d\lambda_d(x)
 =\int_{V_h}(\delta_h^jf)\varphi\,\mathrm d\lambda_d(x).
\]
Rozwinięcie różnicy iloczynu daje prawie wszędzie w $V_h$
\[
 \delta_h^j(A\nabla u)(x)
 =A(x)\nabla w_h(x)
   +(\delta_h^jA)(x)\nabla u(x+he_j).
\]
Po przeniesieniu drugiego składnika uzyskujemy
\begin{equation}\label{eq:pde15-diff-eq}
 \begin{aligned}
 \int_{V_h} A\nabla w_h\cdot\nabla\varphi\,\mathrm d\lambda_d(x)
 &=\int_{V_h}(\delta_h^jf)\varphi\,\mathrm d\lambda_d(x)\\
 &\quad-\int_{V_h}(\delta_h^jA)(x)\nabla u(x+he_j)
       \cdot\nabla\varphi\,\mathrm d\lambda_d(x).
 \end{aligned}
\end{equation}
Dla ustalonego $h$ współczynniki $\delta_h^jA$ są ograniczone,
a $\delta_h^jf$ i $\nabla w_h$ należą do $L^2(V_h)$
(z oszacowania norm obu translacji przez normy na $V$).
Zatem obie strony równania są ciągłe względem normy
$H^1(V_h)$. Na mocy lematu~\ref{lem:pde9-compact-product}
funkcję $\varphi=\eta^2w_h$ można aproksymować testami
gładkimi o nośnikach w ustalonym zwartym podzbiorze $V_h$.
Jej przedłużenie przez zero na $\R^d$ należy do
$H^1(\R^d)$, a słaby gradient jest przedłużeniem przez zero
pola
\[
 \nabla\varphi=\eta^2\nabla w_h+2\eta w_h\nabla\eta.
\]
Wolno więc użyć tej funkcji w~\eqref{eq:pde15-diff-eq}.

\emph{Krok 2: jednostajna kontrola różnic współczynników.}
Niech $D=\|\partial_jA\|_{L^\infty(V;\mathrm{op})}$,
gdzie używamy normy operatorowej macierzy. Dla ustalonego
$h$ mamy $|\delta_h^jA(x)|_{\mathrm{op}}\le D$ prawie
wszędzie na $K$. Istotnie, wygładzamy $A$ przez splot
na otoczeniu zwartego zbioru $K+\overline{B_\rho(0)}$
zawartego w $V$. Dla dostatecznie małego parametru
jądro nieujemne o całce $1$ i komutacja pochodnych
ze splotem (twierdzenie~\ref{thm:pde7-mollify})
dają $\|\partial_j A_\varepsilon(x)\|_{\mathrm{op}}\le D$
na tym otoczeniu. Całkując po odcinku od $x$ do $x+he_j$,
otrzymujemy analogiczne ograniczenie dla
$\delta_h^j A_\varepsilon(x)$. Z twierdzenia
~\ref{thm:pde9-mollify} dla $p=2$ (na relatywnie zwartym
otoczeniu) mamy zbieżność $A_\varepsilon$ do $A$ w $L^2$
na $K$ i $K+he_j$. W konsekwencji
$\delta_h^j A_\varepsilon\to\delta_h^j A$ w $L^2(K)$.
Po przejściu do podciągu zbieżnego prawie wszędzie
otrzymujemy żądane ograniczenie. Stała $D$ nie zależy od $h$.

\emph{Krok 3: oszacowanie energetyczne.}
Oznaczmy
\[
 X=\|\eta\nabla w_h\|_{2,V_h},\quad
 Y=\|\nabla u\|_{2,V},\quad F=\|f\|_{2,V},
 \quad M=\|\nabla\eta\|_\infty .
\]
Lemat~\ref{lem:pde15-difference} i niezmienniczość miary
względem translacji dają
$\|w_h\|_{2,K}\le Y$ oraz
$\|\nabla u(\cdot+he_j)\|_{2,K}
=\|\nabla u\|_{2,K+he_j}\le Y$.
Ponieważ $0\le\eta\le1$,
\[
 \|\nabla\varphi\|_{2,V_h}\le X+2MY.
\]
Nie szacujemy normy $\delta_h^jf$ przez stałą niezależną
od $h$, gdyż $f$ może nie mieć żadnej słabej pochodnej.
Zamiast tego, na mocy dyskretnego całkowania przez części,
\[
 \left|\int_{V_h}(\delta_h^jf)\varphi\,\mathrm d\lambda_d(x)\right|
 =\left|\int_V f\,\delta_{-h}^j\varphi\,\mathrm d\lambda_d(x)\right|
 \le F\|\partial_j\varphi\|_{2,V_h}
 \le F(X+2MY).
\]
W przedostatniej nierówności zastosowaliśmy oszacowanie
ilorazu różnicowego z lematu~\ref{lem:pde15-difference}
do przedłużenia $\varphi$ przez zero na $\R^d$:
biorąc rosnące zwarte kule i przechodząc do granicy,
dostajemy oszacowanie na całej $\R^d$ przez normę
$\partial_j\varphi$, która znika poza $K$.

Rozwijając lewą stronę~\eqref{eq:pde15-diff-eq}
i korzystając z eliptyczności na $V$, otrzymujemy
\[
 \lambda X^2
 \le 2\Lambda MXY+F(X+2MY)+DY(X+2MY).
\]
Pierwszy składnik oszacowano przez Cauchy'ego--Schwarza
i $|A\xi|\le\Lambda|\xi|$, ostatni przez ograniczenie
$\delta_h^j A$ i normę przesuniętego gradientu na $K$.
Nie wykorzystujemy symetrii macierzy $A$.
Kładąc $T=F+(2\Lambda M+D)Y$, otrzymujemy
$\lambda X^2\le TX+2MY(F+DY)$.
Nierówność Younga daje
$TX\le\lambda X^2/2+T^2/(2\lambda)$, zatem
\[
 X^2\le \frac{T^2}{\lambda^2}
       +\frac{4M}{\lambda}Y(F+DY)
 \le C_1(F^2+Y^2).
\]
Ostatnia nierówność wynika z
$(a+b)^2\le2a^2+2b^2$ oraz $2FY\le F^2+Y^2$;
$C_1$ zależy wyłącznie od $\lambda,\Lambda,M,D$.
Wybierając maksimum po skończenie wielu kierunkach $j$,
możemy przyjąć jedną stałą $C_1$ niezależną od $j,h$.

\emph{Krok 4: odzyskanie drugich słabych pochodnych.}
Ponieważ $\eta=1$ na $U$, dla każdego $j,k$ mamy
\[
 \|\delta_h^j\partial_ku\|_{2,U}^2
 =\|\partial_kw_h\|_{2,U}^2\le C_1(F^2+Y^2).
\]
Ustalmy $j,k$ i ciąg niezerowych $h_n\to0$ z $|h_n|<\rho$.
Z ograniczoności w przestrzeni Hilberta $L^2(U)$
wybieramy słabo zbieżny podciąg. Dla każdego
$\zeta\in C_c^\infty(U)$ dyskretne całkowanie przez części
oraz zbieżność $\delta_{-h_n}^j\zeta\to\partial_j\zeta$
w $L^2(U)$ (z nośnikami w $U$ dla dużych $n$)
identyfikują granicę słabą z $\partial_j\partial_ku$
dokładnie jak w dowodzie lematu~\ref{lem:pde15-criterion}.
Półciągłość normy daje
$\|\partial_j\partial_ku\|_{2,U}^2\le C_1(F^2+Y^2)$.
Dotyczy to wszystkich $d^2$ par indeksów, więc $u\in H^2(U)$.
Dodając normy pochodnych rzędów zero i jeden, otrzymujemy
\[
 \|u\|_{H^2(U)}^2
 \le\|u\|_{H^1(V)}^2+d^2C_1
       \bigl(\|f\|_{2,V}^2+\|\nabla u\|_{2,V}^2\bigr).
\]
Po wyciągnięciu pierwiastka dostajemy
~\eqref{eq:pde15-h2-est}. Wybór $U,V$ był dowolny,
więc $u\in H^2_{\mathrm{loc}}(\Omega)$.
\end{proof}

\begin{uwaga}[Zależność od $V$ w oszacowaniu]
Wszystkie przesunięcia użyte do oszacowania normy
na $U$ mieszczą się w $V$. Nie jest potrzebna wypukłość
ani regularność brzegu $V$. Stała zależy między innymi
od normy gradientu wybranej funkcji odcinającej,
która może wzrastać przy zmniejszaniu odstępu
między $\overline U$ i $\R^d\setminus V$.
\end{uwaga}

\section{Oszacowanie Caccioppolego z prawą stroną}
\label{sec:pde15-caccioppoli}
\begin{stwierdzenie}[Wersja niejednorodna]
Niech $r>0$, $B_{2r}(x_0)\Subset\Omega$,
$A\in L^\infty(B_{2r};\R^{d\times d})$ spełnia warunki
eliptyczności i ograniczoności z definicji~\ref{def:pde15-elliptic},
$f\in L^2(B_{2r})$, a $u\in H^1(B_{2r})$ jest słabym
rozwiązaniem $-\diver(A\nabla u)=f$. Istnieje $C=C(d,\lambda,\Lambda)>0$
taka, że dla każdego $k\in\R$
\[
 \int_{B_r}|\nabla u|^2\,\mathrm d\lambda_d(x)
 \le C\left(r^{-2}\int_{B_{2r}}|u-k|^2\,\mathrm d\lambda_d(x)
 +r^2\int_{B_{2r}}|f|^2\,\mathrm d\lambda_d(x)\right).
\]
W szczególności można wziąć $k=0$ lub średnią $u$ po $B_{2r}$.
\end{stwierdzenie}
\begin{proof}
Wszystkie kule mają środek $x_0$. Kładziemy $v=u-k$;
ponieważ $\nabla v=\nabla u$, funkcja $v$ spełnia to samo
równanie słabe. Wybieramy $\eta\in C_c^\infty(B_{2r})$,
$0\le\eta\le1$, równą $1$ na $B_r$ i spełniającą
$|\nabla\eta|\le C_0/r$, jak w dowodzie twierdzenia
~\ref{thm:pde14-caccioppoli}. Stała $C_0$ jest niezależna
od $r$ i $x_0$. Z lematu~\ref{lem:pde9-compact-product}
funkcja $\varphi=\eta^2v$ należy do $H^1_0(B_{2r})$
i może być aproksymowana gładkimi testami w $H^1$.
Ciągłość obu stron równania słabego, jak w definicji
~\ref{def:pde15-elliptic}, pozwala użyć $\varphi$ jako testu.
Rozwijając jej gradient, otrzymujemy
\[
 \int_{B_{2r}}\eta^2 A\nabla u\cdot\nabla u\,\mathrm d\lambda_d(x)
 =\int_{B_{2r}}f\eta^2v\,\mathrm d\lambda_d(x)
 -2\int_{B_{2r}}\eta v A\nabla u\cdot\nabla\eta\,\mathrm d\lambda_d(x).
\]
Oznaczmy $X=\|\eta\nabla u\|_{2,B_{2r}}$,
$Y=\|v\|_{2,B_{2r}}$ i $F=\|f\|_{2,B_{2r}}$.
Eliptyczność, ograniczoność $A$ i nierówność
Cauchy'ego--Schwarza dają
\[
 \lambda X^2\le FY+2\Lambda C_0r^{-1}XY.
\]
W szczególności nie musimy stosować nierówności
Poincarégo do $\eta v$. Z nierówności Younga mamy
\[
 FY\le\tfrac12 r^2 F^2+\tfrac12 r^{-2}Y^2,\qquad
 2\Lambda C_0r^{-1}XY
 \le\tfrac{\lambda}{2}X^2+
       \tfrac{2\Lambda^2 C_0^2}{\lambda r^2}Y^2.
\]
Przenosząc $\lambda X^2/2$ na lewą stronę, otrzymujemy
\[
 X^2\le \lambda^{-1}r^2 F^2+
 \left(\lambda^{-1}+\frac{4\Lambda^2 C_0^2}{\lambda^2}\right)
 r^{-2}Y^2.
\]
Ponieważ $\eta=1$ na $B_r$, daje to tezę z
$C=\lambda^{-1}+4\Lambda^2C_0^2/\lambda^2$.
Średnia $u$ na $B_{2r}$ minimalizuje $Y^2$ względem $k$,
jak wykazano w dowodzie twierdzenia~\ref{thm:pde14-caccioppoli}.
\end{proof}

\section{Regularność wyższych rzędów i bootstrap}
\label{sec:pde15-bootstrap}
\begin{twierdzenie}[Wewnętrzny bootstrap Sobolewa]
\label{thm:pde15-bootstrap}
Niech $m\in\N_0$, $\Omega\subset\R^d$ będzie obszarem,
$A\in W^{m+1,\infty}_{\mathrm{loc}}(\Omega;\R^{d\times d})$
będzie jednostajnie eliptyczna i ograniczona na zwartych
podzbiorach, $f\in H^m_{\mathrm{loc}}(\Omega)$,
a $u\in H^1_{\mathrm{loc}}(\Omega)$ niech spełnia
$-\diver(A\nabla u)=f$ słabo. Wtedy
$u\in H^{m+2}_{\mathrm{loc}}(\Omega)$.
Dla niepustych otwartych zbiorów $U\Subset V\Subset\Omega$
istnieje $C>0$ takie, że
\[
 \|u\|_{H^{m+2}(U)}
 \le C(\|f\|_{H^m(V)}+\|u\|_{H^1(V)}).
\]
Stała zależy od $m,d$, $U,V$, stałych eliptyczności na $V$
i normy $A$ w $W^{m+1,\infty}(V)$.
\end{twierdzenie}
\begin{proof}
Dowodzimy indukcyjnie jednocześnie regularności i oszacowania.
Dla $m=0$ jest to twierdzenie~\ref{thm:pde15-h2}.
Niech $m\ge1$ i załóżmy oba wnioski dla $m-1$.
Ponieważ obecne założenia implikują założenia dla $m-1$,
mamy $u\in H^{m+1}_{\mathrm{loc}}(\Omega)$.

\emph{Krok 1: różniczkowanie dystrybucyjne iloczynu.}
Ustalmy wielowskaźnik $\alpha$, $|\alpha|=m$, i połóżmy
$v_\alpha=\partial^\alpha u\in H^1_{\mathrm{loc}}(\Omega)$.
Uzasadnimy wzór Leibniza dla $A\nabla u$.
Jeżeli $a\in W^{1,\infty}(W)$ oraz $z\in H^1(W)$,
gdzie $W\Subset\Omega$ jest otwarty, to obie funkcje
należą do $H^1(W)$, gdyż $W$ ma skończoną miarę.
Z twierdzenia~\ref{thm:pde9-products} dla $p=q=2$
otrzymujemy słabą regułę iloczynu
$\partial_j(az)=(\partial_j a)z+a\partial_j z$.
Iloczyn i oba składniki pochodnej należą do $L^2(W)$,
ponieważ $a,\partial_j a\in L^\infty(W)$.
Stąd $az\in H^1(W)$ i
\[
 \|az\|_{H^1(W)}
 \le C_d\|a\|_{W^{1,\infty}(W)}\|z\|_{H^1(W)}.
\]
Stosujemy ten argument składowo również do macierzy.
Powtarzając regułę iloczynu, dla $|\gamma|\le m$ dostajemy
dystrybucyjnie
\[
 \partial^\gamma(A\nabla u)=
 \sum_{\beta\le\gamma}\binom{\gamma}{\beta}
 (\partial^\beta A)\nabla\partial^{\gamma-\beta}u .
\]
Istotnie, dla $|\gamma|<m$ każdy składnik ma oba czynniki
o regularności wystarczającej do kolejnego zastosowania
reguły: $\partial^\beta A\in W^{1,\infty}(W)$ i
$\nabla\partial^{\gamma-\beta}u\in H^1(W)$.
Po zróżniczkowaniu względem jednej współrzędnej
współczynniki sumują się według tożsamości Pascala
dla odpowiedniej składowej wielowskaźnika.
Wszystkie potrzebne pochodne istnieją już na mocy
założenia indukcyjnego, nie zakładamy więc z góry
regularności $H^{m+2}$.

\emph{Krok 2: równanie dla pochodnej i regularność.}
Ponieważ pochodne dystrybucyjne są przemienne,
z równania wyjściowego i powyższego wzoru wynika
\[
 -\diver(A\nabla v_\alpha)=\partial^\alpha f+\diver Q_\alpha,
 \qquad
 Q_\alpha=\sum_{0<\beta\le\alpha}\binom{\alpha}{\beta}
 (\partial^\beta A)\nabla\partial^{\alpha-\beta}u .
\]
Dla $0<\beta\le\alpha$ mamy
$\partial^\beta A\in W^{1,\infty}_{\mathrm{loc}}$
oraz $\nabla\partial^{\alpha-\beta}u\in
H^{|\beta|}_{\mathrm{loc}}\subset H^1_{\mathrm{loc}}$.
Z oszacowania iloczynu z kroku 1 wynika dla każdego
$W\Subset\Omega$
\[
 \|Q_\alpha\|_{H^1(W)}
 \le C_{d,m}\|A\|_{W^{m+1,\infty}(W)}\|u\|_{H^{m+1}(W)}.
\]
Dywergencja pola $H^1(W;\R^d)$ należy do $L^2(W)$
i jej norma nie przekracza $\sqrt d$ razy normy
macierzy pierwszych pochodnych w $L^2(W)$
(z nierówności Cauchy'ego--Schwarza dla sumy $d$ składników).
Zatem prawa strona $F_\alpha=\partial^\alpha f+\diver Q_\alpha$
jest lokalnie kwadratowo całkowalna.
Twierdzenie~\ref{thm:pde15-h2} zastosowane do $v_\alpha$
daje $v_\alpha\in H^2_{\mathrm{loc}}(\Omega)$.
Każdy wielowskaźnik długości $m+2$ jest sumą
wielowskaźnika długości $m$ i wielowskaźnika długości $2$,
więc wraz z założeniem indukcyjnym daje to
$u\in H^{m+2}_{\mathrm{loc}}(\Omega)$.

\emph{Krok 3: oszacowanie normy.}
Dla ustalonych $U,V$ wybieramy otwarty zbiór $W$ taki,
że $U\Subset W\Subset V$. Wtedy twierdzenie
~\ref{thm:pde15-h2} oraz oszacowanie $Q_\alpha$ dają
\[
 \|v_\alpha\|_{H^2(U)}
 \le C_1(\|F_\alpha\|_{2,W}+\|v_\alpha\|_{H^1(W)})
 \le C_2(\|f\|_{H^m(W)}+\|u\|_{H^{m+1}(W)}).
\]
Stałe są niezależne od $u,f$; maksimum po skończenie wielu
$\alpha$ można przyjąć jako wspólną stałą.
Z oszacowania indukcyjnego dla pary $W,V$ mamy
\[
 \|u\|_{H^{m+1}(W)}
 \le C_3(\|f\|_{H^{m-1}(V)}+\|u\|_{H^1(V)}).
\]
Normy $f$ na $W$ i rzędu $m-1$ na $V$ nie przekraczają
$\|f\|_{H^m(V)}$. Sumując kwadraty oszacowań
$\|v_\alpha\|_{H^2(U)}$ oraz dodając oszacowanie indukcyjne
dla pochodnych niższych rzędów, otrzymujemy tezę
po wyciągnięciu pierwiastka. Liczba wielowskaźników
zależy tylko od $d,m$; wybór $W$ zależy od $U,V$.
To zamyka indukcję również dla zależności stałych.
\end{proof}

\begin{wniosek}[Gładkość wewnętrzna]
Jeżeli $A\in C^\infty(\Omega;\R^{d\times d})$ jest lokalnie
jednostajnie eliptyczna, $f\in C^\infty(\Omega)$ i
$u\in H^1_{\mathrm{loc}}(\Omega)$ jest słabym rozwiązaniem,
to klasa $u$ ma jedyny gładki reprezentant na $\Omega$.
Reprezentant ten spełnia równanie również klasycznie.
\end{wniosek}
\begin{proof}
Każda pochodna współczynników i prawej strony jest
ograniczona na domknięciu zbioru $V\Subset\Omega$.
Poprzednie twierdzenie stosujemy dla każdego $m$
i otrzymujemy $u\in H^k_{\mathrm{loc}}(\Omega)$
dla wszystkich $k\in\N$.
Ustalmy kulę $B\Subset\Omega$. Dla $d\ge2$ zastosujemy
twierdzenie~\ref{thm:pde10-domain-embed} do wszystkich
słabych pochodnych na $B$. Jeżeli wszystkie należą do
$L^p(B)$ dla pewnego skończonego $p\ge2$, to wszystkie
należą do $W^{1,p}(B)$, bo ich kolejne słabe pochodne
również należą do $L^p(B)$. Gdy $p<d$, możemy zastąpić
$p$ przez $dp/(d-p)$, czyli zmniejszyć $1/p$ o $1/d$.
Po skończenie wielu krokach osiągamy $p\ge d$;
jeśli $p=d$, stosujemy zanurzenie do $L^{d+1}(B)$.
W ten sposób wszystkie pochodne należą do $L^p(B)$
dla pewnego $p>d$, więc $u\in W^{k,p}(B)$ dla każdego $k$.
Wniosek~\ref{cor:pde10-Morrey-domain} daje dla każdego
$\ell\ge0$ reprezentanta klasy $C^\ell$ na $B$.
Dla $d=1$ stosujemy do każdej pochodnej na przedziale $B$
jednowymiarowe twierdzenie o reprezentancie Sobolewa
(w szczególności wniosek~\ref{cor:pde9-trace-oneD}):
każda pochodna jest ciągła, a reprezentant poprzedniej
jest absolutnie ciągły i ma tę pochodną prawie wszędzie.
Podstawowe twierdzenie rachunku całkowego daje wtedy
klasyczną pochodną równą ciągłemu reprezentantowi.
Iteracja daje gładkość. W obu przypadkach ciągłe
reprezentanty tej samej klasy są równe punktowo na
przecięciach otwartych zbiorów, gdyż są równe prawie
wszędzie. Sklejamy je więc do jednego reprezentanta
$C^\infty(\Omega)$. Różnica $-\diver(A\nabla u)-f$
dla tego reprezentanta jest gładka i generuje dystrybucję
zerową; z ciągłości i iniektywności dystrybucji regularnych
(stwierdzenie~\ref{prop:pde7-injective}) znika wszędzie.
\end{proof}

\section{Regularność przy brzegu: metoda spłaszczania}
\label{sec:pde15-boundary}

\begin{lemat}[Forma dywergencyjna po spłaszczeniu brzegu]
\label{lem:pde15-flatten}
Niech $V,W\subset\R^d$ będą zbiorami otwartymi i
$\Phi:W\to V$ będzie dyfeomorfizmem klasy $C^{1,1}$
(tj. klasy $C^1$ z lokalnie lipschitzowską pochodną).
Piszemy $J=D\Phi$ i $q=|\det J|$. Jeżeli
$u\in H^1_{\mathrm{loc}}(V)$ jest słabym rozwiązaniem
$-\diver(A\nabla u)=f$ na $V$, gdzie
$A\in W^{1,\infty}_{\mathrm{loc}}(V;\R^{d\times d})$
i $f\in L^2_{\mathrm{loc}}(V)$, to
$v=u\circ\Phi\in H^1_{\mathrm{loc}}(W)$ i spełnia równanie
$-\diver(B\nabla v)=F$ słabo, gdzie
\[
 B(y)=q(y)J(y)^{-1}A(\Phi(y))J(y)^{-\mathsf T},
 \qquad F(y)=q(y)f(\Phi(y)).
\]
Mamy $B\in W^{1,\infty}_{\mathrm{loc}}(W)$ i
$F\in L^2_{\mathrm{loc}}(W)$.
Jeżeli na mierzalnym $E\subset W$ zachodzą ograniczenia
$\|J\|_{\mathrm{op}}\le M$, $\|J^{-1}\|_{\mathrm{op}}\le N$
i na $\Phi(E)$ macierz $A$ spełnia warunki z definicji
~\ref{def:pde15-elliptic} ze stałymi $\lambda,\Lambda$,
to na $E$ macierz $B$ spełnia te same warunki ze stałymi
$\lambda'= \lambda/(N^d M^2)$ i
$\Lambda'=\Lambda M^d N^2$.
Jeżeli powyższe ograniczenia $M,N$ zachodzą na całym $W$,
to składanie z $\Phi$ jest ograniczonym operatorem
$H^1(V)\to H^1(W)$ oraz $F\in L^2(W)$ dla $f\in L^2(V)$.
Jeżeli dodatkowo $D^2\Phi$ jest słabo ograniczone na $W$,
to składanie z $\Phi$ jest ograniczone również
$H^2(V)\to H^2(W)$. Analogiczne twierdzenia stosują się
do $\Phi^{-1}$ przy odpowiednich ograniczeniach jej pochodnych.
\end{lemat}
\begin{proof}
\emph{Krok 1: słaba reguła łańcuchowa.}
Z klasycznej zmiany zmiennych dla dyfeomorfizmu $C^1$
wynika zachowywanie zbiorów miary zero przez $\Phi$
i $\Phi^{-1}$. Złożenie jest więc dobrze określone dla
klas funkcji równych prawie wszędzie. Niech $W_0\Subset W$
będzie otwarty. Zwartość $\Phi(\overline{W_0})\subset V$
pozwala wybrać otwarty $V_0\Subset V$ zawierający ten obraz.
Na $\overline{W_0}$ macierze $J,J^{-1}$ są ograniczone
i $q$ jest oddzielone od zera. Jeżeli $M,N$ są ich
odpowiednimi ograniczeniami, to z oszacowań wartości
osobliwych mamy $N^{-d}\le q\le M^d$.
Dla każdej funkcji $z\in L^2(V_0)$ wzór zmiany zmiennych daje
\[
 \|z\circ\Phi\|_{2,W_0}^2
 =\int_{\Phi(W_0)}|z(x)|^2
        q(\Phi^{-1}(x))^{-1}\,\mathrm d\lambda_d(x)
 \le N^d\|z\|_{2,V_0}^2.
\]
Oszacowanie stosuje się składowo także do pól wektorowych.
Z twierdzenia~\ref{thm:pde9-mollify} wybieramy funkcje
$u_n$ gładkie na $V_0$, zbieżne do $u$ w $H^1(V_0)$.
Klasyczna reguła łańcuchowa daje
$\nabla(u_n\circ\Phi)=J^{\mathsf T}(\nabla u_n)\circ\Phi$.
Powyższe oszacowanie zapewnia zbieżność złożeń w $L^2(W_0)$
do $v=u\circ\Phi$ i ich gradientów do
$J^{\mathsf T}(\nabla u)\circ\Phi$ w $L^2(W_0)$.
Przechodząc do granicy w tożsamości definiującej słabe
pochodne, otrzymujemy
\[
 \nabla v=J^{\mathsf T}(\nabla u)\circ\Phi,
 \qquad (\nabla u)\circ\Phi=J^{-\mathsf T}\nabla v.
\]
Ponieważ $W_0$ był dowolny, wzory zachodzą lokalnie
na całym $W$. Przy globalnych ograniczeniach $M,N$
ten sam wzór zmiany zmiennych na $W$ daje
\[
 \|v\|_{2,W}\le N^{d/2}\|u\|_{2,V},\qquad
 \|\nabla v\|_{2,W}\le MN^{d/2}\|\nabla u\|_{2,V}.
\]
W szczególności $v\in H^1(W)$ dla $u\in H^1(V)$.

\emph{Krok 2: transformacja równania.}
Dla $\psi\in C_c^\infty(W)$ funkcja
$\varphi=\psi\circ\Phi^{-1}$ jest klasy $C^1$,
ma zwarty nośnik w $V$ i należy do $H^1(V)$.
Wolno jej użyć na mocy rozszerzenia klasy testów
z definicji~\ref{def:pde15-elliptic}. Reguła łańcuchowa
daje $(\nabla\varphi)\circ\Phi=J^{-\mathsf T}\nabla\psi$.
W równaniu słabym zmieniamy zmienne $x=\Phi(y)$:
\[
 \int_W q\bigl(A(\Phi)J^{-\mathsf T}\nabla v\bigr)
      \cdot J^{-\mathsf T}\nabla\psi\,\mathrm d\lambda_d(y)
 =\int_W q(f\circ\Phi)\psi\,\mathrm d\lambda_d(y).
\]
Z tożsamości $a\cdot J^{-\mathsf T}b=(J^{-1}a)\cdot b$
otrzymujemy dokładnie $\int_W B\nabla v\cdot\nabla\psi
\,\mathrm d\lambda_d(y)=\int_W F\psi\,\mathrm d\lambda_d(y)$.
Ponadto przy ograniczeniach $M,N$ na $W_0$ mamy
\[
 \|F\|_{2,W_0}^2
 =\int_{\Phi(W_0)}q(\Phi^{-1}(x))|f(x)|^2\,\mathrm d\lambda_d(x)
 \le M^d\|f\|_{2,\Phi(W_0)}^2.
\]
To daje lokalną kwadratową całkowalność $F$, a dla
globalnego $M$ również oszacowanie na całym $W$.

\emph{Krok 3: eliptyczność i regularność współczynników.}
Na $E$ mamy prawie wszędzie, dla każdego $\xi\in\R^d$,
\[
 \xi^{\mathsf T}B\xi
 =q(J^{-\mathsf T}\xi)^{\mathsf T}A(\Phi)(J^{-\mathsf T}\xi)
 \ge N^{-d}\lambda M^{-2}|\xi|^2,
 \qquad |B\xi|\le M^d N^2\Lambda|\xi|.
\]
W pierwszej nierówności użyto
$|\xi|=|J^{\mathsf T}J^{-\mathsf T}\xi|
\le M|J^{-\mathsf T}\xi|$. Zbiór wyjątkowy dla $A$
po przeniesieniu przez $\Phi^{-1}$ pozostaje zbiorem miary
zero, wspólnym dla wszystkich $\xi$; symetria $A$
nie jest potrzebna.

Lokalnie współczynniki $A$ mają lipschitzowskie
reprezentanty: dla $d\ge2$ stosujemy
twierdzenie~\ref{thm:pde10-domain-embed} na małych kulach
z $p=\infty$, a dla $d=1$ jednowymiarową teorię Sobolewa.
Na przecięciach reprezentanty zgadzają się punktowo,
gdyż są ciągłe i równe prawie wszędzie.
Złożenie $A\circ\Phi$ jest lokalnie lipschitzowskie,
ponieważ $\Phi$ jest lokalnie lipschitzowska jako
odwzorowanie klasy $C^1$. Macierz $J$ jest lokalnie
lipschitzowska z założenia; wyznacznik jest wielomianem
jej współczynników, a wartość bezwzględna jest
lipschitzowska. Macierz odwrotna jest lokalnie
lipschitzowska na mocy tożsamości
\[
 J(y)^{-1}-J(z)^{-1}
 =J(y)^{-1}(J(z)-J(y))J(z)^{-1}
\]
i lokalnej ograniczoności $J^{-1}$.
W konsekwencji wszystkie czynniki we wzorze na $B$
są lokalnie lipschitzowskie. Z twierdzenia Rademachera
i charakterystyki słabych pochodnych funkcji
lipschitzowskich (stosując na odcinkach jednowymiarową
regułę całkowania przez części i twierdzenie Fubiniego)
wynika $B\in W^{1,\infty}_{\mathrm{loc}}(W)$.

\emph{Krok 4: przenoszenie regularności drugiego rzędu.}
Dla $u\in H^2_{\mathrm{loc}}(V)$ stosujemy krok 1
także do każdej składowej $\partial_i u$. Z reguły
iloczynu dla mnożnika $W^{1,\infty}$ udowodnionej
w dowodzie twierdzenia~\ref{thm:pde15-bootstrap} wynika
\[
 \partial_{ab}v
 =\sum_{i,k=1}^d(\partial_{ik}u)\circ\Phi\,
       \partial_a\Phi_i\,\partial_b\Phi_k
   +\sum_{i=1}^d(\partial_i u)\circ\Phi\,\partial_{ab}\Phi_i.
\]
Wszystkie składniki są lokalnie kwadratowo całkowalne,
więc $v\in H^2_{\mathrm{loc}}(W)$.
Jeżeli $M,N$ i $L=\|D^2\Phi\|_\infty$ są skończone
na całym $W$, wzór zmiany zmiennych z kroku 1
i powyższa skończona suma dają
\[
 \|D^2 v\|_{2,W}
 \le C_d N^{d/2}
      \bigl(M^2\|D^2u\|_{2,V}+L\|\nabla u\|_{2,V}\bigr).
\]
Wraz z oszacowaniem $H^1$ dowodzi to ograniczoności
składania w $H^2$. Dla odwrotnej zmiany zmiennych
powtarzamy argument z $\Phi^{-1}$.
\end{proof}

\begin{lemat}[Oszacowanie pochodnych stycznych w półkuli]
\label{lem:pde15-tangential}
Niech $d\ge2$, $B_2^+=\{x\in\R^d:|x|<2,\ x_d>0\}$,
$A\in W^{1,\infty}(B_2^+;\R^{d\times d})$ będzie
jednostajnie eliptyczna, $f\in L^2(B_2^+)$ oraz
$v\in H^1(B_2^+)$ niech spełnia
\[
 \int_{B_2^+}A\nabla v\cdot\nabla\varphi\,\mathrm d\lambda_d(x)
 =\int_{B_2^+}f\varphi\,\mathrm d\lambda_d(x)
 \quad(\varphi\in H_0^1(B_2^+)).
\]
Jeżeli ślad $v$ na płaskiej części brzegu $\{x_d=0,|x|<2\}$
wynosi zero, to dla każdego $j<d$
$\partial_j\nabla v\in L^2(B_1^+)$ i
\[
 \|\partial_j\nabla v\|_{L^2(B_1^+)}
 \le C\bigl(\|f\|_{L^2(B_2^+)}+\|v\|_{H^1(B_2^+)}\bigr),
\]
gdzie $C$ zależy od $d$, eliptyczności oraz
$\|A\|_{W^{1,\infty}(B_2^+)}$.
\end{lemat}
\begin{proof}
Ustalmy $j<d$. Wybieramy $\eta\in C_c^\infty(B_{3/2})$
równą $1$ na $B_1$, i dla małych $|h|$ kładziemy
$w_h=\delta_h^jv$. Ponieważ przesunięcia są styczne do
płaskiego brzegu, $w_h$ ma tam zerowy ślad.
W szczególności $\eta^2w_h\in H_0^1(B_2^+)$.
Ponadto $-\delta_{-h}^j(\eta^2w_h)$ również należy
do $H_0^1(B_2^+)$: jego nośnik omija część sferyczną,
a ślad na płaskiej jest zerowy.
Testujemy słabe równanie ostatnią funkcją i wykonujemy
dyskretne całkowanie przez części. Otrzymujemy
\[
 \int_{B_2^+}A\nabla w_h\cdot\nabla(\eta^2w_h)
 \,\mathrm d\lambda_d(x)
 =\int_{B_2^+}(\delta_h^jf)\eta^2w_h\,\mathrm d\lambda_d(x)
 -\int_{B_2^+}(\delta_h^jA)\nabla v(\cdot+he_j)
 \cdot\nabla(\eta^2w_h)\,\mathrm d\lambda_d(x).
\]
W pierwszym członie po prawej nie różniczkujemy $f$:
\[
 \left|\int_{B_2^+}(\delta_h^jf)\eta^2w_h
 \,\mathrm d\lambda_d\right|
 =\left|\int_{B_2^+}f\,\delta_{-h}^j(\eta^2w_h)
 \,\mathrm d\lambda_d\right|
 \le\|f\|_2\|\partial_j(\eta^2w_h)\|_2.
\]
Własności ilorazów różnicowych dają
$\|w_h\|_{L^2(\operatorname{supp}\eta)}\le\|\nabla v\|_2$
i $\|\delta_h^jA\|_\infty\le\|\partial_jA\|_\infty$.
Rozwijamy $\nabla(\eta^2w_h)=\eta^2\nabla w_h+
2\eta w_h\nabla\eta$, stosujemy eliptyczność do
pierwszego członu i nierówność Younga do pozostałych.
Po pochłonięciu otrzymujemy jednostajnie w $h$
\[
 \|\eta\nabla w_h\|_2^2
 \le C(\|f\|_2^2+\|v\|_{H^1(B_2^+)}^2).
\]
Kryterium z lematu~\ref{lem:pde15-criterion}, stosowane
lokalnie w półkuli i z oszacowaniem jednostajnym aż do
płaskiego brzegu, daje $\partial_j\partial_kv\in L^2(B_1^+)$
dla każdego $k$. Nie rozciągamy funkcji przez zero
w kierunku normalnym, ponieważ takie przesunięcie
nie zachowuje równania słabego.
\end{proof}



\begin{twierdzenie}[Globalne oszacowanie $H^2$ dla Dirichleta]
\label{thm:pde15-boundary}
Niech $\Omega\subset\R^d$ będzie ograniczonym obszarem
klasy $C^{1,1}$, $A\in W^{1,\infty}(\Omega;\R^{d\times d})$
będzie jednostajnie eliptyczna i ograniczona, a
$f\in L^2(\Omega)$. Jeżeli $u\in H_0^1(\Omega)$ spełnia
$-\diver(A\nabla u)=f$ słabo, to $u\in H^2(\Omega)$ oraz
\begin{equation}\label{eq:pde15-boundary}
 \|u\|_{H^2(\Omega)}\le C\bigl(\|f\|_{L^2(\Omega)}
 +\|u\|_{H^1(\Omega)}\bigr),
\end{equation}
gdzie $C$ zależy od $d,\Omega$, stałych eliptyczności
i $\|A\|_{W^{1,\infty}(\Omega)}$.
\end{twierdzenie}
\begin{proof}
Podzielimy dowód na kroki, rozdzielając dwa różne zagadnienia:
kontrolę pochodnych stycznych i odzyskanie pochodnej normalnej.

\emph{Krok 1: skończone pokrycie brzegu.}
Ponieważ $\partial\Omega$ jest zwarty i klasy $C^{1,1}$,
można pokryć go skończenie wieloma otoczeniami, w których
po obrocie i translacji $\Omega$ ma postać
$x_d>\gamma(x')$, gdzie $\gamma\in C^{1,1}$.
Odwzorowanie $y' =x'$, $y_d=x_d-\gamma(x')$
i jego odwrotność są klasy $C^{1,1}$, a ich pierwsze
pochodne są Lipschitzowskie.
Dokładny wzór na nową macierz $B$ i prawą stronę $F$
podaje lemat~\ref{lem:pde15-flatten}; w tej zmianie zmiennych
$\det D\Phi=1$. Otrzymujemy jednostajnie eliptyczne
współczynniki $B\in W^{1,\infty}$ na mniejszym otoczeniu.
Na płaskim brzegu $\{y_d=0\}$ funkcja $v$ ma ślad zerowy.

\emph{Krok 2: pochodne styczne.}
Ustalmy funkcję odcinającą $\eta$ w półkuli $B_R^+$.
Dla $j<d$ przesunięcia $y\mapsto y+he_j$ nie przecinają
płaskiego brzegu. Iloraz różnicowy $w_h=\delta_h^jv$
ma zatem zerowy ślad na tym brzegu. Dopuszczalny jest test
$\varphi=-\delta_{-h}^j(\eta^2w_h)$ w słabym równaniu;
w przypadku $\eta^2w_h$ o zwartym nośniku wewnątrz
półkuli z wyjątkiem płaskiej części brzegu test należy do
$H_0^1(B_R^+)$.
Stosujemy lemat~\ref{lem:pde15-tangential} po przeskalowaniu
półkuli i odcięciu równania w mniejszym otoczeniu.
Otrzymujemy $\partial_j\nabla v\in L^2$ lokalnie do
płaskiego brzegu dla $j<d$ i kontrolę tych pochodnych
przez normy $\|f\|_2+\|v\|_{H^1}$.
Lemat uzasadnia dokładnie, dlaczego wolno przenieść
iloraz różnicowy na test, bez zakładania pochodnych $f$.

\emph{Krok 3: druga pochodna normalna.}
Dla $d=1$ nie ma kierunków stycznych: z równania
$-(a v')'=F$ mamy $(a v')'=-F\in L^2$, a ponieważ
$a\in W^{1,\infty}$ i $a\ge\lambda>0$, wynika
$v''\in L^2$. Dalej rozważamy $d\ge2$.
Rozpiszmy równanie (w zmiennych $y$) jako
\[
 -a_{dd}\partial_{dd}v=f+
 \sum_{(i,j)\ne(d,d)}\partial_i(a_{ij}\partial_jv)
 +(\partial_da_{dd})\partial_dv.
\]
Wszystkie składniki prawej strony oprócz
$\partial_{dd}v$ są już w $L^2$: gdy $i<d$,
korzystamy z pochodnych stycznych $\partial_i\nabla v$,
a gdy $i=d$ i $j<d$, korzystamy z przemienności
słabych pochodnych mieszanych.
Ponadto $a_{dd}\ge\lambda>0$ prawie wszędzie.
Stąd $\partial_{dd}v\in L^2$, a więc $v\in H^2$
lokalnie aż do spłaszczonego brzegu.

\emph{Krok 4: sklejenie oszacowań.}
Wnętrze $\Omega$ pokrywamy skończenie wieloma mniejszymi
zbiorami, a brzeg powyższymi mapami. Wybieramy partycję jedności
i sumujemy lokalne oszacowania. Ponieważ zmiany zmiennych
są klasy $C^{1,1}$, przeniesienie drugich słabych pochodnych
z powrotem do współrzędnych $x$ zachowuje $H^2$.
Daje to~\eqref{eq:pde15-boundary}.
\end{proof}

\begin{uwaga}[Niejednorodne dane Dirichleta]
Jeśli $g$ jest śladem pewnej funkcji $G\in H^2(\Omega)$,
to dla $v=u-G\in H_0^1(\Omega)$ otrzymujemy
$-\diver(A\nabla v)=f+\diver(A\nabla G)\in L^2(\Omega)$.
Stosując twierdzenie~\ref{thm:pde15-boundary} do $v$,
otrzymujemy $u\in H^2(\Omega)$. Zamiast wymagania
$G\in H^2$ można sformułować warunek na $g$ w odpowiedniej
przestrzeni śladów wyższego rzędu; wymaga to osobnego
twierdzenia o operatorze śladu.
\end{uwaga}


\begin{uwaga}[Granica obowiązywania identyfikacji dziedziny]
Dla $A=I$ twierdzenie~\ref{thm:pde15-boundary}
wyjaśnia identyfikację $D(-\Delta_D)=H^2(\Omega)\cap
H_0^1(\Omega)$ na obszarze klasy $C^{1,1}$, zapowiedzianą
w uwadze~\ref{rem:pde12-corners}. Na ogólnym obszarze
Lipschitzowskim taka identyfikacja może być fałszywa;
nie wolno jej używać w argumentach spektralnych
bez odpowiednich dodatkowych założeń.
\end{uwaga}

\section{Dwa kontrprzykłady: współczynniki i geometria brzegu}
\label{sec:pde15-counterexamples}
\begin{kontrprzyklad}[Skok współczynnika: $H^1$ bez $H^2$]
\label{ex:pde15-jump}
Na $\Omega=(-1,1)$ połóżmy
\[
 a(x)=\begin{cases}1,&x<0,\\2,&x\ge0,\end{cases}
 \qquad
 u(x)=\begin{cases}x,&x\le0,\\x/2,&x>0.\end{cases}
\]
Funkcja $u$ jest Lipschitzowska, więc należy do $H^1(\Omega)$.
Mamy $a(x)u'(x)=1$ prawie wszędzie, zatem
$-(a u')'=0$ dystrybucyjnie.
Jednak $u'$ ma skok równy $-1/2$ w zerze, a więc
$u''=-\tfrac12\delta_0$ jako dystrybucja; stąd
$u\notin H^2(\Omega)$. Samo ograniczenie i eliptyczność
$A$ nie gwarantują drugich pochodnych w $L^2$.
\end{kontrprzyklad}

\begin{kontrprzyklad}[Narożnik wklęsły]
\label{ex:pde15-corner}
Niech $\Omega$ będzie ograniczonym obszarem płaskim,
który w sąsiedztwie zera jest sektorem o kącie rozwarcia
$3\pi/2$, np. obszarem w kształcie litery L.
Współrzędne biegunowe dobieramy tak, że w tym sektorze
$0<\theta<3\pi/2$. Funkcja
$w(r,\theta)=r^{2/3}\sin(2\theta/3)$ jest harmoniczna dla $r>0$
i znika na ramionach sektora. Wybieramy gładką radialną
funkcję $\chi$, równą $1$ w otoczeniu zera i zerową poza
małą kulą zawartą w sąsiedztwie narożnika. Wtedy
$u=\chi w\in H_0^1(\Omega)$ oraz
$f=-\Delta u\in L^2(\Omega)$: blisko zera mamy $f=0$,
a tam, gdzie $\nabla\chi\ne0$, wszystkie składniki są gładkie.
Tymczasem $|\nabla^2u|\asymp r^{-4/3}$ blisko zera na
podsektorze dodatniej miary kątowej i dlatego
\[
 \int_{\Omega\cap B_\eps(0)}|\nabla^2u|^2\,\mathrm d\lambda_2
 \ge c\int_0^\eps r^{-8/3}r\,\mathrm dr=+\infty.
\]
Zatem $u\notin H^2(\Omega)$ mimo gładkiego operatora
$-\Delta$ i prawej strony $f\in L^2$. Obszar Lipschitzowski
nie musi mieć globalnej regularności $H^2$.
\end{kontrprzyklad}

\section{Wnioski i zakres dalszych twierdzeń}
\label{sec:pde15-summary}
Regularność $H^2$ we wnętrzu wynika z eliptyczności i
Lipschitzowskich współczynników; nie trzeba zakładać
$C^2$-regularności rozwiązania. Regularność przy brzegu
wymaga natomiast dodatkowej geometrii brzegu i poprawnie
postawionych warunków brzegowych.
W kolejnym rozdziale przejdziemy od oszacowań $L^2$ do
regularności $W^{2,p}$ oraz oszacowań Schaudera.
Nie należy przy tym mylić regularności Sobolewa z analitycznością:
$C^\infty$-gładkość współczynników gwarantuje gładkość rozwiązania
w powyższej sytuacji, lecz sama nie gwarantuje analityczności.

